\section{Proposed Method}
In this section, we present the proposed Cross-Modality Semantic Alignment Transformer (CMSA-Trans) for zero-shot fault diagnosis of rotating machinery. The core idea is to bridge the gap between vibration signal space and technical fault semantics generated by Large Language Models (LLMs), enabling the model to recognize unseen fault categories without labeled vibration samples.

\subsection{Overall Architecture}
The CMSA-Trans architecture is designed to map vibration features and text-based semantic prototypes into a unified high-dimensional latent space. As illustrated in the conceptual framework, the system consists of three primary modules: a Signal Encoder, a Semantic Prototyping module, and a Cross-Modality Alignment (CMA) module. The process begins with raw vibration signals $X \in \mathbb{R}^{B \times L}$, where $B$ is the batch size and $L$ is the signal length. These signals are partitioned into non-overlapping temporal segments and projected into a latent sequence by the Signal Encoder, which leverages a Time-Series Transformer to capture complex long-range dependencies inherent in mechanical wave patterns. Simultaneously, technical descriptions of various fault modes (e.g., "inner race spalling under high-load conditions") are processed via an LLM-based Semantic Prototyping module to extract representative semantic vectors $V_s$. The CMA module then utilizes a mutual attention mechanism to align the rhythmic features of the vibration signal with the descriptive attributes of the fault semantics. By minimizing the distance between the signal embedding and its corresponding semantic prototype, the network learns a modality-invariant representation. Finally, a zero-shot classification head performs inference by calculating the cosine similarity between the test signal embedding and the pre-computed semantic prototypes of all candidate classes, allowing for the identification of faults that were not present during the training phase. This architecture effectively transforms fault diagnosis from a closed-set classification problem into a semantic-driven retrieval task. The overall framework is built to accommodate various mechanical components, including bearings, gears, and shafts, by adapting the semantic descriptions accordingly. The scalability of the Transformer-based backbone ensures that the model can handle multi-channel sensor inputs, further enhancing its diagnostic precision in complex industrial environments. By integrating LLM-derived knowledge, CMSA-Trans mitigates the reliance on expensive labeled datasets, marking a shift toward more autonomous and adaptable health monitoring systems.

\subsection{Signal Encoder: Time-Series Transformer}
The vibration signals generated by rotating machinery are non-stationary and contain transient impulses that reflect the health state of components like bearings and gears. To effectively extract discriminative features, we employ a Time-Series Transformer-based encoder. Given a raw input signal sequence $x = [x_1, x_2, \dots, x_L]$, we first apply a patch-based embedding layer. The signal is reshaped into $N$ patches of length $P$, such that $L = N \times P$. Each patch is then linearly projected into a vector of dimension $D$. To preserve the temporal order of mechanical events, a learnable positional encoding $E_{pos} \in \mathbb{R}^{N \times D}$ is added to the patch embeddings. The resulting sequence of input embeddings $Z_0$ is formulated as:
\begin{equation}
Z_0 = [x_p^1 W_E; x_p^2 W_E; \dots; x_p^N W_E] + E_{pos}
\end{equation}
where $W_E \in \mathbb{R}^{P \times D}$ is the trainable projection matrix. The sequence $Z_0$ is passed through a stack of $M$ Transformer layers. Each layer consists of a multi-head self-attention (MHSA) block followed by a position-wise feed-forward network (FFN). Layer normalization (LN) and residual connections are applied before and after each block. The self-attention mechanism allows the encoder to weigh the importance of different temporal segments, capturing the periodic characteristics of vibration signals that are often obscured by noise. The output of the final Transformer layer provides a robust high-level representation of the signal, denoted as $F_{sig} \in \mathbb{R}^{N \times D}$. By utilizing the global receptive field of the Transformer, the encoder can capture long-term correlations across multiple shaft rotations, which is critical for diagnosing complex mechanical faults in unsteady operating conditions \cite{vaswani2017attention, li2023transformer}. Furthermore, the Time-Series Transformer architecture is specifically adapted to handle the high sampling rates typically encountered in vibration monitoring. Unlike conventional convolutional neural networks (CNNs), which may struggle with very long sequences due to limited kernel sizes, the self-attention mechanism enables the direct modeling of dependencies between distant signal segments, such as the recurring impacts of a rolling element passing over a localized race defect. This makes the encoder particularly adept at distinguishing between health states with similar spectral profiles but different temporal signatures.

\subsection{Semantic Prototyping via LLM}
A key challenge in zero-shot fault diagnosis is the lack of vibration data for novel fault modes. We address this by leveraging the extensive knowledge encoded in LLMs to generate technical semantic prototypes. For each target fault category $C_i$, we construct a technical prompt that describes the physical manifestations and mechanical characteristics associated with the fault (e.g., "characteristic frequency components of rolling element defect in a tapered roller bearing with localized pitting"). The LLM produces a detailed technical description $T_i$, which is then converted into a fixed-length semantic vector (prototype) $p_i$ using a pre-trained sentence embedding model or the text encoder of a CLIP-like architecture. This process ensures that the semantic space is populated with vectors that represent the underlying physics of the mechanical failure rather than just simple label names. The set of semantic prototypes $P = \{p_1, p_2, \dots, p_K\}$ for $K$ classes is frozen during training to maintain a stable target for alignment. This approach bypasses the need for high-frequency vibration data collection for every possible failure scenario, instead relying on the linguistic descriptions of mechanical phenomena \cite{radford2021learning, brown2020language}. The extracted prototypes $V_s$ serve as the "anchor" points in the latent space, toward which the signal embeddings are guided during the training process. The sophistication of the LLM-generated descriptions is crucial; they must capture not only the fault type but also the influence of operating parameters such as rotational speed and load, which modify the vibration signature. By enriching the semantic prototypes with these physical insights, we ensure that the alignment process is grounded in the actual mechanical behavior of the equipment. This methodology represents a significant departure from traditional zero-shot methods that rely on hand-crafted attributes, providing a more scalable and expertise-driven framework for modern Industry 4.0 applications.

\subsection{Cross-Modality Alignment}
The CMA module is the bridge that aligns the signal features $F_{sig}$ with the semantic prototypes $V_s$. We implement a cross-attention mechanism where the signal features act as the Query ($Q$), and the semantic prototypes act as the Key ($K$) and Value ($V$). Specifically, we project $F_{sig}$ into the semantic space using a projection layer $W_P \in \mathbb{R}^{D \times d_s}$, where $d_s$ is the dimensionality of the semantic prototypes. The cross-attention scores are calculated as follows:
\begin{equation}
Q_{sig} = F_{sig} W_Q , \quad K_{sem} = V_s W_K , \quad V_{sem} = V_s W_V
\end{equation}
The attention output is defined as:
\begin{equation}
\text{Attention}(Q_{sig}, K_{sem}, V_{sem}) = \text{softmax}\left(\frac{Q_{sig} K_{sem}^T}{\sqrt{d_k}}\right) V_{sem}
\end{equation}
where $d_k$ is the scaling factor. To further refine the alignment, we maximize the cosine similarity between the global signal embedding $z_{sig}$ (obtained by global average pooling over the attended sequence) and the corresponding ground-truth semantic prototype $p_i$:
\begin{equation}
S(z_{sig}, p_i) = \frac{z_{sig} \cdot p_i}{\|z_{sig}\| \|p_i\|}
\end{equation}
Through the cross-modality attention, the model learns to selectively attend to the vibration features that are most consistent with the technical descriptions of fault modes. This mechanism enables the "semantic grounding" of mechanical signals, allowing the network to recognize signal patterns that align with the high-level descriptions of unseen fault classes during inference \cite{li2024crossmodality}. The alignment process is critical, as it transforms the discrete diagnostic problem into a continuous representation learning task. By projecting the features into a shared semantic space, we enable the model to generalize based on the conceptual similarities between different fault conditions. For instance, a bearing defect and a gear tooth breakage might share certain impulsive characteristics, and the semantic prototypes help the model learn to distinguish between them by focusing on the unique descriptive elements provided by the LLM. This focus on cross-modal consistency ensures that the learned signal embeddings are not just discriminative of the training set but are truly representative of the underlying mechanical health states.

\subsection{Training and Optimization}
The CMSA-Trans is trained using a cross-entropy-based contrastive loss, which encourages the network to map signals near their correct semantic prototypes while pushing them away from irrelevant ones. For a batch of vibration signals and their corresponding semantics, the loss function is formulated as:
\begin{equation}
\mathcal{L} = - \sum_{j=1}^{B} \log \frac{\exp(S(z_{sig}^j, p_y^j) / \tau)}{\sum_{k=1}^{K} \exp(S(z_{sig}^j, p_k) / \tau)}
\end{equation}
where $\tau$ is a learnable temperature parameter and $y$ is the ground-truth label index. During optimization, the weights of the Signal Encoder and the CMA projection layers are updated via backpropagation using an Adam optimizer. By optimizing Eq. (7), the model develops a robust joint embedding space where signal samples of the same fault category cluster around their respective LLM-generated semantic prototype. This enables the zero-shot capability: at test time, for an input signal of an unseen class $C_{unseen}$, the model computes its embedding and assigns the label of the nearest prototype in the semantic space. The learning rate is typically scheduled with a cosine decay to ensure stable convergence. This strategy ensures that the diagnostic capability generalizes to categories outside the training set, provided their technical semantics are defined \cite{sohn2016improved}. This training paradigm shift from traditional classification to semantic alignment allows for dynamic updates to the fault dictionary without retraining the entire signal encoder, offering immense flexibility for evolving industrial systems. Every optimization step reinforces the link between the temporal dynamics of the vibration signal and the technical concepts of machinery maintenance, resulting in a model that is both high-performing and physically interpretable.